\section{System Description}
\label{sec:system}

The experimental platform is \system{}, an open-source framework that exposes
both vanilla RAG and KG-RAG under a shared interface.
The two systems use the same corpora, embeddings, and generation model; they
differ only in retrieval and evidence organisation.

\begin{figure*}[t]
\centering
\begin{tikzpicture}[
  font=\footnotesize\sffamily,
  >={Stealth[length=2.2mm,width=1.7mm]},
  stage/.style={draw=black!50, line width=0.5pt, rounded corners=2pt,
                fill=white, minimum height=0.95cm, minimum width=1.9cm,
                align=center, inner xsep=6pt, inner ysep=4pt},
  kgstage/.style={stage, draw=teal!55!black, line width=0.7pt, fill=teal!7},
  flow/.style={->, line width=0.55pt, black!55},
  rail/.style={teal!55!black, line width=0.55pt},
  rowlab/.style={font=\scriptsize\sffamily, text=black!45},
]

%% ── Build row ────────────────────────────────────────────────────────────────
\node[stage]   (corpus)  at (0,2.55)   {Domain\\corpus};
\node[stage]   (chunk)   at (3.2,2.55) {Passage\\chunking};
\node[stage]   (extract) at (6.4,2.55) {Entity \& relation\\extraction};
\node[kgstage] (kg)      at (9.6,2.55) {Knowledge\\graph};

%% ── Query row ────────────────────────────────────────────────────────────────
\node[stage] (q)    at (0,0)     {Question $q$};
\node[stage] (link) at (2.55,0)  {Entity\\linking};
\node[stage] (sub)  at (5.1,0)   {Subgraph\\retrieval};
\node[stage] (ctx)  at (7.65,0)  {Context\\assembly};
\node[stage] (llm)  at (10.2,0)  {LLM\\generation};
\node[stage] (ans)  at (12.6,0)  {Answer $\hat{a}$};

%% ── Row labels ───────────────────────────────────────────────────────────────
\node[rowlab, rotate=90, anchor=center] at (-1.55,2.55) {build time};
\node[rowlab, rotate=90, anchor=center] at (-1.55,0)    {query time};

%% ── Flow arrows ──────────────────────────────────────────────────────────────
\draw[flow] (corpus) -- (chunk);
\draw[flow] (chunk)  -- (extract);
\draw[flow] (extract) -- (kg);
\draw[flow] (q)    -- (link);
\draw[flow] (link) -- (sub);
\draw[flow] (sub)  -- (ctx);
\draw[flow] (ctx)  -- (llm);
\draw[flow] (llm)  -- (ans);

%% ── Graph-state rail: the KG conditions three query-time stages ──────────────
\coordinate (raily) at (0,1.25);
\draw[rail] (kg.south) -- (kg.south |- raily);
\draw[rail] (link.north |- raily) -- (kg.south |- raily);
\draw[->, rail] (link.north |- raily) -- (link.north);
\draw[->, rail] (sub.north  |- raily) -- (sub.north);
\draw[->, rail] (ctx.north  |- raily) -- (ctx.north);
\node[font=\scriptsize\itshape\sffamily, text=teal!45!black,
      above, inner sep=1.5pt] at (8.62,1.25) {graph state};

%% ── Sampling loop: N samples re-enter retrieval, returning ~the same state ──
\draw[dashed, line width=0.55pt, black!45, ->]
  (llm.south) .. controls +(0,-0.85) and +(0,-0.85) .. (link.south);
\node[font=\scriptsize\itshape\sffamily, text=black!50, fill=white,
      inner sep=1.5pt] at (6.4,-1.07)
  {$\times N$ samples --- entity-first routing returns nearly the same subgraph};

%% ── Measurement taps: where each diagnostic family attaches ─────────────────
\tikzset{tap/.style={draw, line width=0.5pt, rounded corners=2pt,
                     align=center, font=\scriptsize\sffamily,
                     inner xsep=4pt, inner ysep=2.5pt, fill=white},
         probe/.style={densely dashed, line width=0.6pt, ->}}
\node[tap, draw=teal!60!black, text=teal!40!black] (tgps) at (5.1,-2.05)
  {GPS\\(retrieval-state)};
\node[tap, draw=violet!60!black, text=violet!40!black] (tseu) at (7.65,-2.05)
  {SEU\\(evidence-state)};
\node[tap, draw=blue!50!black, text=blue!35!black] (tans) at (12.0,-2.05)
  {DSE, SD-UQ, VN-Ent.\\(answer-state, over $N$ samples)};
\draw[probe, teal!60!black]   (sub.south) -- (tgps.north);
\draw[probe, violet!60!black] (ctx.south) -- (tseu.north);
\draw[probe, blue!50!black]   (ans.south) -- (tans.north);

\end{tikzpicture}
\caption{%
  \textbf{\system{} pipeline with the three measurement taps.}
  \emph{Build time} (top): a domain corpus is chunked and converted via
  ontology-guided extraction into a typed knowledge graph with passage-level
  provenance.
  \emph{Query time} (bottom): the question is linked to seed entities, and the
  graph state conditions retrieval at three points---entity linking, subgraph
  traversal, and context assembly---before the LLM produces the final answer.
  The dashed probes mark where each diagnostic family attaches: GPS scores
  graph support in risk orientation, SEU scores the assembled evidence against
  the answer, and the answer-state estimators see only the $N$ sampled answers.
  The sampling loop is the lock-in mechanism: each of the $N$ samples re-enters
  retrieval, and entity-first routing returns nearly the same subgraph, so
  answer-state uncertainty mostly reflects decoder noise under almost-fixed
  evidence (Section~\ref{sec:metrics}).%
}
\label{fig:pipeline}
\end{figure*}


\subsection{\system{} Overview}
\label{sec:ontographrag_overview}

\system{} is designed around a simple principle: graph retrieval should add
structure without discarding dense evidence (Figure~\ref{fig:pipeline}).
The graph makes retrieval commitments explicit: recognised entities, accepted
relations, triples, paths, and licensing passages.
During KG construction, passages are converted into entity and relation nodes
with ontology-guided typing, anchor grounding, confidence scores, provenance
links, and contradiction flags when a schema is available.
During retrieval, the system first tries to route the question through matched
entities and local graph neighbourhoods.
When the anchor is weak, it uses dense retrieval and retriever-first graph
expansion as fallback.
The resulting context contains both textual passages and explicit graph paths,
so the answer can be evaluated at three levels: the answer state, the
evidence state, and the retrieval state.
\system{} logs anchors, paths, and supporting passages with the final answer,
so the provenance trace is part of the retrieval output, not a private prompt
construction detail.

\subsection{Vanilla RAG}
\label{sec:vanilla_rag}

Vanilla RAG performs direct vector retrieval over chunk embeddings.
For each question, it retrieves the top-$k$ chunks above a similarity
threshold and optionally appends adjacent chunks to capture answers split
across chunk boundaries.
The LLM therefore receives flat text only.
It has no explicit entities, relations, or multi-hop paths, and thus no
graph-side object on which to define retrieval-state diagnostics.

\subsection{KG-RAG}
\label{sec:kg_rag}

KG-RAG replaces flat retrieval with a routed hybrid pipeline rather than a
single graph walk.

\paragraph{Stage 1: seed selection.}
The preferred path is entity-first anchoring: the system identifies seed
entities with symbolic matching and embedding lookup over short mentions.
If no reliable anchor is found, the system does not immediately reduce to
vector-only retrieval; instead it can route to a retriever-first graph pass.

\paragraph{Stage 2: local graph expansion and scoring.}
Starting from the seed entities, the system traverses the graph up to a
dataset-specific hop limit, producing an explicit retrieval state of linked
chunks, entities, relations, and readable traversal paths.
Traversal is provenance-aware: on question-bundle datasets, paths and edges are
restricted to the current question's local evidence rather than allowed to
borrow support from other questions in the same dataset graph.
Chunks are scored by hop distance and local diffusion over the retrieved
entity subgraph.

\paragraph{Stage 3: retriever-first graph expansion when anchoring is weak.}
When entity-first retrieval is weak or empty, the system first retrieves dense
text chunks, extracts their linked entities, expands the graph from those
passage-derived seeds, and re-scores the resulting context with a combined
vector-plus-graph signal.

\paragraph{Stage 4: evidence organisation.}
The final prompt is not a flat mixture of chunks and graph paths.
Instead, retrieved graph paths are grouped with overlapping supporting passages
into explicit reasoning chains, with remaining passages listed separately as
additional evidence.
The generator sees both structured multi-hop support and the local text that
grounds each chain.

\paragraph{Diagnostic role.}
\system{} is useful here because it makes the retrieval state explicit rather
than tacit.
The graph is not treated as an oracle; it is an auditable substrate whose
triples and paths may be helpful or wrong, but are diagnostically richer than
an unstructured list of nearest-neighbour chunks.
Because the system logs which entities were matched, which paths were
traversed, and which route was taken, the retrieval state is an explicit,
queryable object rather than an internal variable.

\paragraph{Retrieval stabilisation.}
The diagnostic question arises when the retrieval state is stable across
repeated calls; retriever-first expansion, vector fallback, and iterative
decomposition can weaken that stabilisation when the graph anchor is brittle,
so the experiments compare diagnostics across retrieval regimes.
Entity-first retrieval can be especially prone to stabilising: once the system
commits to an anchor entity, the subsequent graph expansion is largely
determined, so the same neighbourhood, right or wrong, tends to recur across
samples.
